\documentclass[11pt,a4paper]{article}

\usepackage[utf8]{inputenc}
\usepackage{amsmath,amssymb,amsthm}
\usepackage{bm}
\usepackage{graphicx}
\usepackage{booktabs}
\usepackage{hyperref}
\usepackage{geometry}
\geometry{margin=1in}

\hypersetup{
    colorlinks=true,
    linkcolor=blue,
    citecolor=blue,
    urlcolor=blue
}

\title{\textbf{Mathematical Predictions and Experimental Design for In-Context Benign Overfitting}}
\author{\textbf{Research Team}}
\date{\today}

\begin{document}

\maketitle

\begin{abstract}
This document outlines the mathematical framework and experimental designs for studying \textit{in-context benign overfitting} in in-context linear regression under a feature-selection mechanism. We analyze the asymptotic behavior of in-distribution (ID) and out-of-distribution (OOD) generalization risks in the high-dimensional proportional scaling regime. Additionally, we specify a series of controlled experiments to validate theoretical predictions, sweep signal-to-noise ratios (SNR), map training dynamics, and investigate OOD task shifts such as Markovian sequences.
\end{abstract}

\section{Introduction}
In-context learning (ICL) allows pre-trained models to solve downstream tasks at inference time by conditioning on a sequence of demonstration examples without updating their weights. The generalization behavior of ICL displays a fascinating dichotomy:
\begin{itemize}
    \item \textbf{Retrieval Mode (In-Distribution / ID Generalization):} Generalizing to prompts generated from tasks that were present during pre-training. Success here indicates that the model has memorized the training task set and retrieves the correct task mapping.
    \item \textbf{Learning Mode (Out-of-Distribution / OOD Generalization):} Generalizing to prompts generated from entirely novel tasks unseen during pre-training. Success here shows the model acts as a meta-learner, learning the task mechanism from the prompt itself.
\end{itemize}

By analogy to classical supervised learning, \textit{in-context benign overfitting} occurs when an overparameterized model memorizes pre-training tasks (achieving zero or near-zero ID risk) while simultaneously generalizing well to novel tasks (achieving low OOD risk).

This document serves as the formal experimental and mathematical design for validating these properties, sweeping parameters like model capacity, SNR, and tracking optimization dynamics.

---

\section{Mathematical Framework}

\subsection{Data Generation Model}
We study in-context linear regression. Let the input sequence for a prompt be:
\begin{equation}
    X = [x_1, y_1, x_2, y_2, \dots, x_T, y_T, x_q]
\end{equation}
where $T$ is the context length (number of demonstrations), and $x_q := x_{T+1}$ is the query vector. 
\begin{itemize}
    \item Covariates are isotropic Gaussian: $x_t \sim \mathcal{N}(0, \Sigma_x)$ with $\Sigma_x = \sigma_x^2 I_d$.
    \item Target labels are generated by $y_t = w^\top x_t + \epsilon_t$, where $\epsilon_t \sim \mathcal{N}(0, \sigma_n^2)$ is label noise.
    \item The true task parameter is $w \sim D_{\text{task}} = \mathcal{N}(0, \Sigma_w)$ with $\Sigma_w = \sigma_w^2 I_d$.
\end{itemize}
Following standard scaling, we fix $\sigma_x = 1/\sqrt{d}$ and $\sigma_w = 1$. The Signal-to-Noise Ratio (SNR) is governed by the relation:
\begin{equation}
    \text{SNR} = \frac{\text{Var}(w^\top x)}{\sigma_n^2} = \frac{\sigma_w^2 \sigma_x^2 d}{\sigma_n^2} = \frac{1}{\sigma_n^2}
\end{equation}

\subsection{Feature Selection \& Capacity Control}
To represent model capacity variations, we employ a feature-selection mechanism. The learner restricts its input to a subset of indices $S \subset [d]$ of size $p = |S|$, chosen uniformly without replacement.
The model creates a feature embedding:
\begin{equation}
    h_S(X) = \operatorname{vec}\left( x_q[S], \hat{w}_{\text{avg}}[S]^\top \right) \in \mathbb{R}^{p^2}
\end{equation}
where $\hat{w}_{\text{avg}}[S] = \frac{1}{T} \sum_{t=1}^T y_t x_t[S]$ is the empirical averaging estimator restricted to $S$. The prediction is linear in these embedded features:
\begin{equation}
    f_{\theta, S}(X) = \theta^\top h_S(X)
\end{equation}
where $\theta \in \mathbb{R}^{p^2}$ represents the trainable parameters of the model.

\subsection{Optimization Regimes}
Given a pre-training dataset of $N$ sequences generated from a pool of $M$ pre-training tasks $\{w_i\}_{i=1}^M \sim D_{\text{task}}$, we minimize the empirical risk:
\begin{equation}
    \mathcal{L}_{\text{emp}}(\theta) = \frac{1}{N} \| H_S \theta - y \|^2
\end{equation}
where $H_S \in \mathbb{R}^{N \times p^2}$ contains the feature embeddings $h_S(X_i)$.
\begin{itemize}
    \item \textbf{Underparameterized Regime ($N > p^2$):} The solution is the unique least-squares estimator:
    \begin{equation}
        \hat{\theta}_S = \left( H_S^\top H_S \right)^{-1} H_S^\top y
    \end{equation}
    \item \textbf{Overparameterized Regime ($N < p^2$):} Gradient descent initialized at zero converges to the minimum-norm interpolating solution:
    \begin{equation}
        \hat{\theta}_S = H_S^\top \left( H_S H_S^\top \right)^{-1} y
    \end{equation}
\end{itemize}

---

\section{Asymptotic Generalization Predictions}
We analyze these estimators in the joint high-dimensional limit where $N, T, M, d, p \to \infty$ at proportional rates:
\begin{equation}
    \frac{N}{d^2} \to \nu, \quad \frac{M}{d} \to \mu, \quad \frac{T}{d} \to \tau, \quad \frac{p}{d} \to \rho
\end{equation}
Under isotropic covariance conditions, we characterize the asymptotic ID and OOD risks. Let $c := \frac{\sigma_n^2 + 1}{\tau}$.

\subsection{Overparameterized Asymptotics ($\nu < \rho^2$)}
In this regime, the empirical risk is interpolated to zero. The limits of the risks are given by:
\begin{equation}
    \mathcal{L}_{\text{OOD}}(\hat{\theta}_S) \xrightarrow{P} \left(1 - \frac{\nu}{\rho^2}\right) \rho (1 + a c) \left[ 1 - c(1 + a + a c) \right] + (1 + a + a c)(\sigma_n^2 + \rho c + 1 - \rho)
\end{equation}
\begin{equation}
    \mathcal{L}_{\text{ID}}(\hat{\theta}_S) \xrightarrow{P} \nu a^2 \bar{\beta}_*^2
\end{equation}
where $a$ and $\bar{\beta}_*$ are the unique positive solutions to the coupled systems derived from the Convex Gaussian Min-Max Theorem (CGMT).

\subsection{Underparameterized Asymptotics ($\nu > \rho^2$)}
In this regime, the features are under-parameterized compared to the number of samples. The asymptotic risks converge to:
\begin{equation}
    \mathcal{L}_{\text{OOD}}(\hat{\theta}_S) \xrightarrow{P} m_c \left[ \kappa_{\infty}^2 (1 + c) - 2 c^2 \rho \right] + (1 + c) \rho c^2 m'_c + \sigma_n^2 + 1 - \rho + \rho c
\end{equation}
\begin{equation}
    \mathcal{L}_{\text{ID}}(\hat{\theta}_S) \xrightarrow{P} \frac{\nu}{\rho^2} \kappa_{\infty}^2
\end{equation}
where $m_c = m_{\gamma}(z_c)$ is the Stieltjes transform of the Marchenko-Pastur law with parameter $\gamma = \rho/\mu$ evaluated at $z_c = -c$, and $m'_c$ is its derivative.

---

\section{Experimental Design}
We present four key experiments to validate the mathematical predictions and characterize the boundaries of benign overfitting.

\subsection{Experiment 1: Model Capacity Sweep (Double Descent)}
\begin{itemize}
    \item \textbf{Objective:} Observe the transition of ID and OOD generalization errors as a function of the relative model scale $\rho^2/\nu$.
    \item \textbf{Setup:} Keep the number of pre-training sequences $N$ fixed, and sweep the model capacity parameter $p$ (and thus $\rho = p/d$) from the underparameterized regime ($\rho^2/\nu < 1$) to the overparameterized regime ($\rho^2/\nu > 1$).
    \item \textbf{Expected Outcome:} A sharp peak (double descent peak) at the interpolation threshold $\rho^2/\nu = 1$. In the overparameterized region, OOD risk should decrease and level off, while ID risk converges to a low value, confirming that scaling leads to benign task memorization.
\end{itemize}

\subsection{Experiment 2: Noise Variance ($\sigma_n^2$) and SNR Sweep}
\begin{itemize}
    \item \textbf{Objective:} Study how label noise affects the generalization limits and determine the transition from benign to harmful overfitting.
    \item \textbf{Setup:} Fix the model scale $\rho^2/\nu > 1$. Sweep the pre-training noise variance $\sigma_n^2$ across a wide range (e.g., $0.1$ to $5.0$). Measure the gap between $\mathcal{L}_{\text{ID}}$ and $\mathcal{L}_{\text{OOD}}$ at convergence.
    \item \textbf{Expected Outcome:} Below a critical noise variance threshold, the overfitting remains benign (generalization holds). Above this threshold, the minimum-norm interpolator starts memorizing pure noise at the cost of OOD representation, leading to a sharp rise in test loss (harmful overfitting).
\end{itemize}

\subsection{Experiment 3: Training Dynamics and Harmful Onset Step}
\begin{itemize}
    \item \textbf{Objective:} Track the evolution of ID and OOD risks over optimization steps during SGD/gradient descent.
    \item \textbf{Setup:} Save periodic model checkpoints during pre-training. Evaluate ID loss and OOD loss at each checkpoint. Scan a grid of evaluation noise scales.
    \item \textbf{Metrics:}
    \begin{itemize}
        \item $\textbf{Best Step:}$ The training step where OOD loss is minimized.
        \item $\textbf{Harmful Onset Step:}$ The training step where the OOD test loss starts to diverge or significantly increase while the train loss continues to decrease.
    \end{itemize}
    \item \textbf{Expected Outcome:} Early training stages primarily learn the clean linear signals. Late stages memorize the residual training-set noise, leading to the onset of harmful overfitting on test sets with noise mismatch.
\end{itemize}

\subsection{Experiment 4: Out-of-Distribution Data Shift (Markovian Inputs)}
\begin{itemize}
    \item \textbf{Objective:} Examine whether the benign overfitting characteristics persist when covariates transition from isotropic Gaussian to sequential Markovian processes.
    \item \textbf{Setup:} Generate prompts where covariates $x_t$ are drawn from a Markov Chain or non-stationary drift processes:
    \begin{equation}
        x_{t+1} = \alpha x_t + \sqrt{1 - \alpha^2} \eta_t, \quad \eta_t \sim \mathcal{N}(0, I_d)
    \end{equation}
    where $\alpha$ is the drift/transition coefficient. Sweep $\alpha$ from $0.0$ (stationary Gaussian) to $0.95$ (highly correlated Markovian).
    \item \textbf{Expected Outcome:} Strong correlation $\alpha$ shifts the effective eigenvalue spectrum of features, altering the interpolation threshold and risk curves. This will determine the robustness limits of in-context learning.
\end{itemize}

---

\section{Parameters \& Simulation Setup}

To maintain reproducibility, we establish the baseline parameter set described in Table~\ref{tab:hyperparameters}.

\begin{table}[h]
\centering
\caption{Hyperparameter Configuration for Benign Overfitting Experiments}
\label{tab:hyperparameters}
\begin{tabular}{lll}
\toprule
\textbf{Category} & \textbf{Parameter} & \textbf{Value / Sweep Range} \\
\midrule
Model Architecture & Ambient Dimension ($d$) & 20 \\
& Attention heads & 2 \\
& Layers & 3 \\
& Embedding dimension & 64 \\
\midrule
Data Generation & Context length ($T$) & 40 (so $T+1 = 81$ tokens) \\
& Noise Standard Deviation ($\sigma_n$) & Sweep: $\{0.1, 0.3, 0.5, 1.0, 1.5, 2.0, 3.0, 5.0\}$ \\
& Task Distribution $D_{\text{task}}$ & $\mathcal{N}(0, I_d)$ \\
& Input Distribution $D_x$ & $\mathcal{N}(0, \frac{1}{d} I_d)$ \\
\midrule
Training Setup & Optimizer & AdamW \\
& Learning Rate & $10^{-4}$ \\
& Batch Size ($B$) & 256 \\
& Training Steps & 500,000 \\
& Checkpoint Frequency & Every 5,000 steps \\
& Random Seeds & $[0, 1, 2, 3, 4]$ \\
\bottomrule
\end{tabular}
\end{table}

\end{document}
